\documentclass[12pt]{article}
\usepackage[a4paper,margin=1in]{geometry}
\usepackage[utf8]{inputenc}
\usepackage[T1]{fontenc}
\usepackage{amsmath,amssymb}
\usepackage{textcomp}
\usepackage{parskip}
\begin{document}
\begin{center}
{\Large \textbf{2009 STEP I}}
\end{center}

\subsection*{Page 1}
\begin{flushleft}
\bigskip\textbf{Section A: Pure Mathematics}\par
\bigskip\textbf{Question 1.} A proper factor of an integer N is a positive integer, not 1 or N , that divides N .\par
2 3\par
\textbf{(i)} Show that 3 5 has exactly 10 proper factors. Determine how many other integers of\par
\[\times\]
the form 3m 5n (where m and n are integers) have exactly 10 proper factors.\par
\[\times\]
\textbf{(ii)} Let N be the smallest positive integer that has exactly 426 proper factors. Determine N ,\par
giving your answer in terms of its prime factors.\par
\bigskip\textbf{Question 2.} A curve has the equation\par
3 3 3 3\par
\[y = x + a + b ,\]
where a and b are positive constants. Show that the tangent to the curve at the point ( a, b) is\par
-\par
2 2 3 3\par
\[b y a x = a + b .\]
-\par
\[In the case a = 1 and b = 2, show that the x-coordinates of the points where the tangent\]
meets the curve satisfy\par
3 2\par
\[7x 3x 27x 17 = 0 .\]
- - -\par
Hence find positive integers p, q, r and s such that\par
3 3 3 3\par
\[p = q + r + s .\]
2 1\par
\bigskip\textbf{Question 3.} (i) By considering the equation x + x a = 0 , show that the equation x = (a x) 2 has\par
- -\par
\[one real solution when a > 0 and no real solutions when a < 0 .\]
Find the number of distinct real solutions of the equation\par
\[x = (1 + a)x a 3\]
-\par
¡ ¢\par
in the cases that arise according to the value of a.\par
\textbf{(ii)} Find the number of distinct real solutions of the equation\par
\[x = (b + x) 2\]
in the cases that arise according to the value of b .\par
\end{flushleft}
\subsection*{Page 2}
\begin{flushleft}
\bigskip\textbf{Question 4.} The sides of a triangle have lengths p q, p and p + q, where p > q > 0 . The largest and\par
-\par
smallest angles of the triangle are α and β, respectively. Show by means of the cosine rule\par
that\par
\[4(1 cos α)(1 cos β) = cos α + cos β .\]
- -\par
In the case α = 2β, show that cos β = and hence find the ratio of the lengths of the sides of\par
the triangle.\par
\bigskip\textbf{Question 5.} A right circular cone has base radius r, height h and slant length `. Its volume V , and the\par
area A of its curved surface, are given by\par
1 2\par
\[V = \pir h , A = \pir` .\]
\textbf{(i)} Given that A is fixed and r is chosen so that V is at its stationary value, show that\par
\[A 2 = 3\pi 2 r 4 and that ` = \sqrt\{3\} r.\]
\textbf{(ii)} Given, instead, that V is fixed and r is chosen so that A is at its stationary value, find h\par
in terms of r.\par
\bigskip\textbf{Question 6.} (i) Show that, for m > 0 ,\par
m x 2 (m 1) 3 (m + 1)\par
\[dx = - + ln m .\]
\[Z 1/m x + 1 2m2\]
\textbf{(ii)} Show by means of a substitution that\par
m 1 m un-1\par
\[dx = du .\]
Z 1/m xn(x + 1) Z 1/m u + 1\par
\textbf{(iii)} Evaluate:\par
2 5\par
x + 3\par
(a) dx ;\par
\[Z 1/2 x^3(x + 1)\]
2 5 3\par
x + x + 1\par
(b) dx .\par
\[Z 1 x^3(x + 1)\]
\end{flushleft}
\subsection*{Page 3}
\begin{flushleft}
\bigskip\textbf{Question 7.} Show that, for any integer m,\par
\[2\pi\]
\[ex cos mx dx = e 2\pi 1 .\]
\[Z 0 m^2 + 1 ¡ - ¢\]
\textbf{(i)} Expand cos(A + B) + cos(A B). Hence show that\par
-\par
\[2\pi\]
\[ex cos x cos 6x dx = 19 e 2\pi 1 .\]
Z 0 650 ¡ - ¢\par
\[2\pi\]
\textbf{(ii)} Evaluate ex sin 2x sin 4x cos x dx .\par
\[\int\]
\bigskip\textbf{Question 8.} (i) The equation of the circle C is\par
2 2 2\par
\[(x 2t) + (y t) = t ,\]
- -\par
\[where t is a positive number. Show that C touches the line y = 0 .\]
Let α be the acute angle between the x-axis and the line joining the origin to the centre\par
\[of C. Show that tan 2α = and deduce that C touches the line 3y = 4x .\]
\textbf{(ii)} Find the equation of the incircle of the triangle formed by the lines y = 0, 3y = 4x and\par
\[4y + 3x = 15 .\]
Note: The incircle of a triangle is the circle, lying totally inside the triangle, that\par
touches all three sides.\par
\end{flushleft}
\subsection*{Page 4}
\begin{flushleft}
\bigskip\textbf{Section B: Mechanics}\par
\bigskip\textbf{Question 9.} Two particles P and Q are projected simultaneously from points O and D, respectively,\par
where D is a distance d directly above O. The initial speed of P is V and its angle of\par
projection above the horizontal is α. The initial speed of Q is kV , where k > 1, and its angle\par
of projection below the horizontal is β. The particles collide at time T after projection.\par
\[Show that cos α = k cos β and that T satisfies the equation\]
2 2 2 2\par
\[(k 1)V T + 2dV T sin α d = 0 .\]
- -\par
Given that the particles collide when P reaches its maximum height, find an expression\par
for sin α in terms of g, d, k and V , and deduce that\par
gd 6 (1 + k)V 2 .\par
\bigskip\textbf{Question 10.} A triangular wedge is fixed to a horizontal surface. The base angles of the wedge are α and\par
\pi α. Two particles, of masses M and m, lie on different faces of the wedge, and are connected\par
-\par
by a light inextensible string which passes over a smooth pulley at the apex of the wedge, as\par
shown in the diagram. The contacts between the particles and the wedge are smooth.\par
m\par
M\par
\[α \pi 2 α\]
-\par
m\par
\textbf{(i)} Show that if tan α > the particle of mass M will slide down the face of the wedge.\par
M\par
2m\par
\textbf{(ii)} Given that tan α = , show that the magnitude of the acceleration of the particles is\par
M\par
g sin α\par
tan α + 2\par
3 3\par
\[and that this is maximised at 4m = M .\]
\end{flushleft}
\subsection*{Page 5}
\begin{flushleft}
\bigskip\textbf{Question 11.} Two particles move on a smooth horizontal table and collide. The masses of the particles\par
are m and M . Their velocities before the collision are ui and vi , respectively, where i is a\par
unit vector and u > v. Their velocities after the collision are pi and qi , respectively. The\par
\[coefficient of restitution between the two particles is e, where e < 1.\]
\textbf{(i)} Show that the loss of kinetic energy due to the collision is\par
m(u p)(u v)(1 e) ,\par
- - -\par
\[and deduce that u > p.\]
\textbf{(ii)} Given that each particle loses the same (non-zero) amount of kinetic energy in the\par
collision, show that\par
\[u + v + p + q = 0 ,\]
\[and that, if m = M ,\]
(M + 3m)u + (3M + m)v\par
\[e = .\]
(M m)(u v)\par
- -\par
\end{flushleft}
\subsection*{Page 6}
\begin{flushleft}
\bigskip\textbf{Section C: Probability and Statistics}\par
\bigskip\textbf{Question 12.} Prove that, for any real numbers x and y, x 2 + y 2 > 2xy .\par
\textbf{(i)} Carol has two bags of sweets. The first bag contains a red sweets and b blue sweets,\par
whereas the second bag contains b red sweets and a blue sweets, where a and b are\par
positive integers. Carol shakes the bags and picks one sweet from each bag without\par
looking. Prove that the probability that the sweets are of the same colour cannot\par
exceed the probability that they are of different colours.\par
\textbf{(ii)} Simon has three bags of sweets. The first bag contains a red sweets, b white sweets and\par
c yellow sweets, where a, b and c are positive integers. The second bag contains b red\par
sweets, c white sweets and a yellow sweets. The third bag contains c red sweets, a white\par
sweets and b yellow sweets. Simon shakes the bags and picks one sweet from each bag\par
without looking. Show that the probability that exactly two of the sweets are of the\par
same colour is\par
2 2 2 2 2 2\par
3(a b + b c + c a + ab + bc + ca )\par
,\par
(a + b + c)3\par
and find the probability that the sweets are all of the same colour. Deduce that the\par
probability that exactly two of the sweets are of the same colour is at least 6 times the\par
probability that the sweets are all of the same colour.\par
\bigskip\textbf{Question 13.} I seat n boys and 3 girls in a line at random, so that each order of the n + 3 children is as\par
likely to occur as any other. Let K be the maximum number of consecutive girls in the line\par
\[so, for example, K = 1 if there is at least one boy between each pair of girls.\]
\textbf{(i)} Find P(K = 3).\par
\textbf{(ii)} Show that\par
n(n 1)\par
\[P(K = 1) = - .\]
(n + 2)(n + 3)\par
\textbf{(iii)} Find E(K).\par
\end{flushleft}
\end{document}
